\documentclass[11pt, a4paper, oneside]{article}
\usepackage{worksheet}
\usepackage[ngerman]{babel}

\begin{document}
	\author{L. Bung}
	\title{OOP: Vererbung (UML) \hspace{10cm} Lösungsvorschlag}
	\subject{SAE}
	\class{E2FI}
	\maketitle
	
	\partnertask{Produktverwaltung eines Gartencenters}
	
	Ein Gartencenter verwaltet ihre Produkte mithilfe eines Softwaresystems, welches überarbeitet werden soll.
	Das System arbeitet aktuell mit den Klassen \textbf{Pflanze}, \textbf{Blumentopf} und \textbf{Werkzeug}.
	
	Sowohl Pflanzen, Blumentöpfe als auch Werkzeuge verfügen über eine Produktnummer, eine Bezeichnung und einen Preis.
	Außerdem kann jedes Produkt dieser Kategorien verkauft werden.
	
	a) Ergänzen Sie die UML-Klassendiagramme um die oben beschriebenen Informationen.
	Fügen Sie zusätzlich sinnvolle Attribute und Methoden zu den einzelnen Klassen hinzu.
	
	\begin{figure}[h]
		\centering
		\begin{tikzpicture}
			\begin{class}[text width=4cm]{Pflanze}{0,0}
				\attribute{produktnummer: int}
				\attribute{bezeichnung: str}
				\attribute{preis: float}
				\attribute{groesse: int}
				\operation{verkaufen()}
				\operation{giessen()}
			\end{class}
			\begin{class}[text width=4cm]{Blumentopf}{5,0}
				\attribute{produktnummer: int}
				\attribute{bezeichnung: str}
				\attribute{preis: float}
				\attribute{volumen: int}
				\operation{verkaufen()}
				\operation{fuellen()}
				\operation{ausleeren()}
			\end{class}
			\begin{class}[text width=4cm]{Werkzeug}{10,0}
				\attribute{produktnummer: int}
				\attribute{bezeichnung: str}
				\attribute{preis: float}
				\attribute{ist\_elektrisch: bool}
				\operation{verkaufen()}
				\operation{benutzen()}
			\end{class}
		\end{tikzpicture}
	\end{figure}
	
	b) Welche Probleme stellen Sie fest? Haben Sie Ideen, wie man diese lösen könnte?
	
	Die Attribute \texttt{produktnummer}, \texttt{bezeichnung} und \texttt{preis} sowie die Methode \texttt{verkaufe()} kommen in jeder der drei Klassen vor, was zu Redundanz führt.
	
	Eine Lösung wäre eine Superklasse (z. B. \texttt{Produkt}), welche die genannten Attribute bzw. Methoden hat und diese an die Subklassen weitergibt.
	
	\pagebreak
	
	\singletask{Erweiterung des Gartencenter-Systems}
	
	Zusätzlich zu den verkauften Produkten bietet die Gärtnerei auch noch \textbf{Dienstleistungen} (zum Beispiel Baumbeschneidungen) und \textbf{Veranstaltungen} (zum Beispiel Kurse zur Gartenpflege) an.
	
	Entwerfen Sie ein UML-Klassendiagramm\footnote{z.B. wieder mithilfe von \url{https://draw.io}} mit mindestens einer Oberklasse und mindestens zwei Unterklassen, bei denen Vererbung sinnvoll eingesetzt wird.
	Überlegen Sie sich sinnvolle Attribute und Methoden für jede Klasse.
	
	\begin{figure}[H]
		\centering
		\begin{tikzpicture}
			\begin{class}[text width=4cm]{Leistung}{5,7}
				\attribute{bezeichnung: str}
				\attribute{preis: float}
			\end{class}
			\begin{class}[text width=4cm]{Dienstleistung}{0,4}
				\inherit{Leistung}
				\attribute{ausfuehrender: str}
				\attribute{datum: date}
				\operation{buchen()}
			\end{class}
			\begin{class}[text width=4cm]{Veranstaltung}{10,4}
				\inherit{Leistung}
				\attribute{leiter: str}
				\attribute{datum: date}
				\operation{reservieren(kunde: str)}
				\operation{stornieren(kunde: str)}
			\end{class}
			\begin{class}[text width=4cm]{Produkt}{5,4}
				\inherit{Leistung}
				\attribute{produktnummer: int}
				\attribute{bezeichnung: str}
				\attribute{preis: float}
				\operation{verkaufen()}
			\end{class}
			\begin{class}[text width=4cm]{Pflanze}{0,0}
				\inherit{Produkt}
				\attribute{groesse: int}
				\operation{giesse()}
			\end{class}
			\begin{class}[text width=4cm]{Blumentopf}{5,0}
				\inherit{Produkt}
				\attribute{volumen: int}
				\operation{fuellen()}
				\operation{ausleeren()}
			\end{class}
			\begin{class}[text width=4cm]{Werkzeug}{10,0}
				\inherit{Produkt}
				\attribute{ist\_elektrisch: bool}
				\operation{benutze()}
			\end{class}
		\end{tikzpicture}
	\end{figure}
	
	\partnertask{Diskussion und Reflexion}
	
	a) Welche Vorteile erkennen Sie beim Einsatz von Vererbung?
	
	\begin{itemize}
		\item \textbf{Weniger Redundanz}: Code wird besser wartbar und weniger fehleranfällig.
		\item \textbf{Strukturierung}: Durch hierarchische Ordnung der Klassen werden reale Situationen besser abgebildet und der Code besser strukturiert.
		\item \textbf{Erweiterbarkeit}: Neue Klassen können leicht hinzugefügt werden.
	\end{itemize}
	
	b) Können Sie sich Fälle vorstellen, in es zu Problemen kommen könnte?
	
	Es kann zu Problemen kommen, wenn eine Klasse von mehreren Klassen erben müsste.
	
\end{document}
